\section{CVP Comparison And Kernel Ambiguity}

CVP asks for a closest vector in an explicitly presented lattice. HCP asks for
a harmonic cochain in the operative Hodge kernel, and in the \HE{} model that
kernel is not explicitly presented by public data.

\subsection{Known-Lattice CVP Surface}

An explicit-lattice CVP instance consists of a generator matrix \(B\) and a
target vector \(t\). The presented lattice is
\[
  \Lambda_B=\{Bx:x\in R^n\}.
\]
A correct CVP answer is a vector \(v\in\Lambda_B\) such that
\[
  \dist(t,v)\le \dist(t,w)
  \qquad
  \text{for every }w\in\Lambda_B.
\]
This is the separation between the known-lattice interface and the
hidden-operator interface: a CVP oracle receives \(B\) and \(t\), while a
public HCP procedure sees only the declared handles and verifier-facing
artifacts unless extra hidden-operator information is supplied.

\subsection{Reduction Direction}

The formalization separates three ingredients. The first is pure linear
algebra: every presented lattice \(\Lambda_B=\image(B)\) can be realized as
the kernel of some matrix \(\Delta\). The second is the presented-operator
Hodge contract: if a solver is sound for the displayed operator \(\Delta\),
then its output is harmonic for that \(\Delta\), hence lies in
\(\ker(\Delta)=\Lambda_B\). The third is the closestness contract owned by the
CVP interface: only after adding minimization among harmonic vectors does the
membership statement become a closest-vector statement. Lean keeps these
surfaces separate.

\begin{proposition}[Kernel realization for explicit lattices]\label{prop:hcp-to-cvp}
  Let \(B\) be an explicit generator matrix over a field and set
  \(\Lambda_B=\image(B)\). There exists a matrix \(\Delta\) such that
  \[
    \ker(\Delta)=\Lambda_B .
  \]
  Consequently, any solver whose output is \(\Delta\)-harmonic for every
  presented \(\Delta\) returns a vector in the explicit lattice
  \(\Lambda_B\) when run on this realized operator. If that solver is also
  closest among \(\Delta\)-harmonic vectors, then its output is a CVP solution
  for \((B,t)\).
\end{proposition}

\begin{proof}
  Let \(B\) be an explicit generator matrix for a lattice
  \[
    \Lambda_B=\{Bx:x\in R^n\}\subseteq R^n
  \]
  and let \(t\in R^n\) be the CVP target. We first show the following
  linear-algebra fact: there is a matrix \(\Delta\) with
  \[
    \ker(\Delta)=\image(B)=\Lambda_B .
  \]
  Put \(V=\image(B)\). Since \(R^n\) is finite-dimensional over a field, choose
  a complementary subspace \(W\) with
  \[
    V+W=R^n,
    \qquad
    V\cap W=0 .
  \]
  Thus every \(y\in R^n\) has a unique decomposition
  \[
    y=v+w,
    \qquad
    v\in V,\quad w\in W.
  \]
  Equivalently, there is a unique element \(p(y)\in W\) such that
  \(y-p(y)\in V\). Define
  \[
    P(y)=p(y).
  \]
  The uniqueness of the \(V\oplus W\) decomposition gives linearity: for
  \(y,z\in R^n\), the vector \(P(y)+P(z)\) lies in \(W\), and
  \[
    (y+z)-(P(y)+P(z))
    =
    (y-P(y))+(z-P(z))\in V,
  \]
  so \(P(y+z)=P(y)+P(z)\). Similarly, for \(a\in R\),
  \(aP(y)\in W\) and
  \[
    ay-aP(y)=a(y-P(y))\in V,
  \]
  so \(P(ay)=aP(y)\).

  We now compute the kernel. If \(y\in V\), then \(0\in W\) and \(y-0\in V\),
  hence uniqueness gives \(P(y)=0\). Conversely, if \(P(y)=0\), then
  \(y-P(y)=y\in V\). Therefore
  \[
    \ker P=V=\image(B).
  \]
  Let \(\Delta\) be the matrix of the linear map \(P\). Then, for every
  \(y\in R^n\),
  \[
    \Delta y=0
    \quad\Longleftrightarrow\quad
    y\in\ker P
    \quad\Longleftrightarrow\quad
    y\in\image(B)
    \quad\Longleftrightarrow\quad
    \exists x,\ Bx=y .
  \]
  This is the content of the Lean lemma
  \lean{exists_matrix_ker_eq_range}.

  Now apply a presented-operator Hodge solver to \(t\) under this realized
  operator. In Lean the solver is a function \(S:R^n\to R^n\), and the
  presented-operator soundness hypothesis is
  \[
    \forall \Delta\, b,\qquad
    \Delta(S(b))=0.
  \]
  Applying this with the \(\Delta\) constructed above and \(b=t\) gives
  \[
    \Delta(S(t))=0.
  \]
  By the kernel calculation, there exists \(x\) with
  \[
    Bx=S(t).
  \]
  Hence \(S(t)\in\Lambda_B\). This is exactly the theorem
  \lean{Hodge_Hardness_Reduction_Proof}: a solver that is harmonic for every
  presented operator yields a lattice vector for every explicit-lattice
  instance once the lattice has been realized as a kernel.

  Finally, suppose in addition that the solver satisfies the closest-harmonic
  contract
  \[
    \forall h',\quad
    \Delta h'=0
    \Longrightarrow
    \dist(t,S(t))\le \dist(t,h').
  \]
  Since the \(\Delta\)-harmonic vectors are exactly the vectors of
  \(\Lambda_B\), this inequality is precisely
  \[
    \dist(t,S(t))\le \dist(t,w)
    \qquad
    \text{for every }w\in\Lambda_B.
  \]
  Together with \(S(t)\in\Lambda_B\), this is the CVP correctness predicate
  \lean{IsClosestLatticeVector} for the explicit instance \((B,t)\).
\end{proof}

Once the hidden operator data has been fixed or presented, the closest harmonic
cochain problem is precisely a CVP-style problem on \(\ker(\Delta)\). Thus CVP
is a known-lattice shadow of the Hodge recovery problem.

\begin{proposition}[Projection-model Hodge recovery is CVP-shaped]\label{prop:hodge-projection-is-cvp}
  Fix a hidden harmonic subspace \(W\subseteq R^n\) with orthogonal projection
  \(P_W\), and let \(\Delta_W=I-P_W\). A cochain \(h\) is
  \(\Delta_W\)-harmonic exactly when \(h\in W\). Therefore, if \(h\) is
  harmonic and closest to \(b\) among all harmonic cochains, then \(h\) is the
  CVP solution for the target \(b\) in the lattice/subspace \(W\).
\end{proposition}

\begin{proof}
  In the projection model,
  \[
    \Delta_W b=(I-P_W)b,
  \]
  equivalently the projection of \(b\) to \(W^\perp\). Hence
  \[
    \Delta_W h=0
    \quad\Longleftrightarrow\quad
    h\in W.
  \]
  The remaining statement is definitional: the closest harmonic representative
  is precisely a vector that is \(\Delta\)-harmonic and minimizes
  \(\dist(b,-)\) among all \(\Delta\)-harmonic vectors.
\end{proof}

\subsection{Oracle Insufficiency}

The strictness direction uses a different contract. A public Hodge solver sees
the cochain-facing observation and may call an explicit-lattice CVP oracle, but
it is not given the hidden operator. Soundness asks that the solver's output be
harmonic for every operative hidden \(\Delta\). Exactness on already-harmonic
inputs asks that, if the observation \(b\) is already harmonic for the operative
\(\Delta\), the solver return \(b\) unchanged.

\begin{proposition}[Known-lattice CVP oracle insufficiency]\label{prop:cvp-insufficient}
  In the hidden-operator model, an oracle for explicit known-lattice CVP does
  not by itself determine the operative Hodge--Laplacian kernel. Consequently,
  such an oracle does not uniformly solve HCP from public data alone.
\end{proposition}

\begin{proof}
  Assume \(\dim R^n\ge 2\), and choose an index \(i\). Let
  \[
    b=e_i,\qquad
    \Delta_1=0,\qquad
    \Delta_2=\operatorname{diag}(1,0,\ldots,0),
  \]
  where the nonzero diagonal entry of \(\Delta_2\) is placed at \(i\). Then
  \[
    \Delta_1 b=0,
    \qquad
    \Delta_2 b=b\ne 0.
  \]
  Thus \(b\) is harmonic for \(\Delta_1\) but not harmonic for
  \(\Delta_2\). If the public projection is lossy enough to identify
  \(\Delta_1\) and \(\Delta_2\), the same cochain-facing observation \(b\) is compatible
  with two hidden Hodge operators that disagree about whether \(b\) lies in the
  harmonic kernel.

  Now suppose, for contradiction, that there were a public-data-only Hodge
  solver \(S\) which may call a correct explicit-lattice CVP oracle \(O\), and
  which is uniformly sound and exact on already-harmonic inputs. Since \(b\) is
  harmonic for \(\Delta_1\), exactness forces
  \[
    S(O,b)=b.
  \]
  But soundness for the hidden operator \(\Delta_2\) forces
  \[
    \Delta_2 S(O,b)=0.
  \]
  Combining the two equalities gives \(\Delta_2 b=0\), contradicting the
  construction of \(b\). Hence no such public-data-only solver exists, even
  assuming the availability of a correct oracle for explicitly presented CVP
  instances.
\end{proof}

The obstruction is kernel ambiguity: the same public observation may be
compatible with multiple hidden operators whose harmonic kernels differ. A CVP
oracle can solve a closest-vector problem after the lattice is presented; it
does not select which hidden kernel public data was meant to present.

\begin{proposition}[Concrete public ambiguity]\label{prop:constant-public-ambiguity}
  In every dimension at least two, the constant public projection is
  public-compatible Hodge ambiguous, hence kernel-ambiguous and lossy.
\end{proposition}

\begin{proof}
  Use the same witnesses as in the proof of
  \cref{prop:cvp-insufficient}: \(b=e_i\), \(\Delta_1=0\), and
  \(\Delta_2=\operatorname{diag}(1,0,\ldots,0)\). The constant projection sends
  every hidden operator to the same public value, so \(\Delta_1\) and
  \(\Delta_2\) are publicly equivalent. Since \(b\) is harmonic for
  \(\Delta_1\) and not for \(\Delta_2\), this gives public-compatible Hodge
  ambiguity. The implications of \cref{prop:ambiguity-hierarchy} then give
  kernel ambiguity and lossiness.
\end{proof}

\subsection{Structural Comparison}

The comparison separates three surfaces:
\begin{enumerate}[label=(\roman*)]
  \item known-lattice closest-vector search;
  \item hidden-operator HCP;
  \item public projection/proximity verification.
\end{enumerate}

The first two surfaces are the hardness story. The third is a verifier-facing
theorem: given accepted hidden/geometric data, later sections construct finite
public shadows and an RS/MCA endpoint without claiming that the projection is
injective.
